\chapter{Git history}
\label{chap:git}
The repository history shows iterative development across data
preparation, model evaluation, application interfaces, and report
documentation. Recent commits updated the API implementation and
\code{.gitignore}, refined report support material, and expanded
project documentation. The current project contains a React
dashboard, a FastAPI service, a Python scoring core, a Streamlit
alternative, a modelling notebook, evaluation artifacts, and the
synthetic data release.

\begin{table}[H]
\centering
\tablefont
\caption{Selected repository milestones}
\label{tab:git-history}
\begin{tabularx}{\textwidth}{|p{0.19\textwidth}|p{0.19\textwidth}|Y|}
\hline
\textbf{Commit} & \textbf{Date} & \textbf{Recorded milestone}\\\hline
\code{e9bb2636} & 2026-10-25 & Updated ignore rules and removed a deprecated supplementary report file.\\\hline
\code{ed6e861e} & 2026-10-25 & Removed deprecated API implementation details.\\\hline
\code{1ebf6156} & 2026-10-23 & Added auxiliary report files and improved document structure.\\\hline
\code{5f4ab2f7} & 2026-10-23 & Refactored report documentation and added supporting details.\\\hline
\code{42d14c99} & 2026-10-22 & Expanded feature dictionary, lineage notes, and report sections.\\\hline
\end{tabularx}
\end{table}

The history indicates that the repository changed during report
preparation. The commit list is a selected snapshot, not a complete
software-development process assessment. A production project should
also use reviewed pull requests, versioned model artifacts, dependency
locks, automated tests, and release tags. The report source in this
file is a new report document; it does not restore or overwrite the
older report source that may exist in repository history.

\chapter{Conclusions}
\label{chap:conclusion}

\section{Summary of the project}
This project developed a local-first prototype for turning several
employee activity logs into a review queue. The data pipeline cleans
and aggregates events into employee-day counts. A notebook labels
synthetic incident-window records, creates user-disjoint development
and test partitions, compares four model families, selects a candidate
and queue capacity on validation data, and persists a model bundle.
The API and dashboard use the saved feature contract to score a new
batch and expose aggregated details to an analyst.

The saved evaluation supports a limited empirical conclusion:
Random Forest achieved the strongest listed performance among the
four models at the selected one-percent test queue capacity on this
CERT r4.2 partition. The result is not a general guarantee and should
not be interpreted as detection accuracy for a real organisation.
The evaluation's positive class is defined by overlap with synthetic
incident windows; both false positives and false negatives are
benchmark-relative.

\section{Learning outcomes}
The implementation provided experience in dataset inspection,
chunked processing, feature engineering, rare-class evaluation,
grouped cross-validation, model persistence, API design, browser
interaction, visual communication, and responsible-use analysis.
It also demonstrated the importance of keeping a training contract
consistent with inference code and separating validation selection
from final test reporting.

\section{Final conclusion}
Forensio demonstrates how a carefully bounded analytics pipeline can
help organise the first stage of a review. Its most important design
property is not the production of a score, but the explicit separation
between statistical prioritisation and human judgment. The current
prototype needs additional software, privacy, security, validation,
and operational controls before any use beyond controlled research
and demonstration.

\chapter{Future scope}
\label{chap:future}
Future work should be driven by evidence and governance needs rather
than by increasing model complexity alone. The following roadmap
describes improvements that would make the system more reliable,
interpretable, and useful while preserving the human-review boundary.

\begin{longtable}{|p{0.12\textwidth}|p{0.28\textwidth}|p{0.42\textwidth}|p{0.12\textwidth}|}
\caption{Prioritised future work}\label{tab:future-work}\\
\hline
\textbf{Phase} & \textbf{Work item} & \textbf{Purpose and acceptance evidence} & \textbf{Priority}\\\hline
\endfirsthead
\hline
\textbf{Phase} & \textbf{Work item} & \textbf{Purpose and acceptance evidence} & \textbf{Priority}\\\hline
\endhead
\hline
\multicolumn{4}{r}{Continued on next page}\\
\endfoot
\hline
\endlastfoot
1 & Add automated unit and integration tests. & Validate parsing, feature alignment, score ordering, API errors, upload staging, and status logic on controlled fixtures. & High\\\hline
1 & Add input schema and quality reports. & Report row counts, invalid fields, duplicate handling, missing sources, and date coverage before scoring. & High\\\hline
1 & Pin dependencies and version model bundles. & Reproduce training and inference with documented library versions, data checksums, and metadata. & High\\\hline
2 & Add model explanation and analyst feedback. & Show feature contributions cautiously and record review outcomes without treating feedback as an automatic label. & Medium\\\hline
2 & Evaluate temporal and cross-release generalisation. & Test robustness across time windows and benchmark versions using predefined protocols. & High\\\hline
2 & Add role and schedule context under governance. & Reduce avoidable confounding only when lawful, necessary, validated, and stakeholder-approved. & Medium\\\hline
3 & Build case-management and audit integration. & Provide authorised review history, reviewer notes, evidence references, and retention controls. & Medium\\\hline
3 & Assess privacy and security controls. & Complete threat modelling, privacy impact review, access control, and deployment testing. & High\\\hline
\end{longtable}

\section{External validation and responsible deployment}
Before operational use, a data owner should establish the lawful and
organisational purpose for analysis and identify who may initiate,
review, and close a case. The organisation should assess privacy,
employment, and security impacts with appropriate stakeholders. It
should define prohibited uses, retention, access, escalation, and a
process for challenging or correcting data. These decisions should be
made before a model score is shown to a decision maker, not inferred
from a dashboard label.

Technical validation should use data that reflect the intended
population, systems, work schedules, and time period. Synthetic CERT
results are valuable for repeatability but are not a substitute for
local validation. A prospective or carefully controlled retrospective
study should define the target outcome, preserve an untouched evaluation
set, measure queue workload, and include uncertainty. It should test
whether apparent performance changes when source coverage, time
boundaries, and user groups vary. If labels are based on past reviews,
the study should consider selection bias and the possibility that
unreviewed activity was never labelled.

Human factors require separate evaluation. Reviewers should understand
that the rank is relative, the labels are incomplete, and model
agreement does not establish intent. The interface should make the
supporting counts and data-quality caveats visible at the point of use.
A reviewer should be able to record a disagreement, request more
context, or close a lead without an adverse outcome. Governance should
monitor whether the system changes how people are treated and whether
review practice remains proportionate to the original purpose.

Operational acceptance should be staged. First, verify technical
reproducibility on a controlled dataset. Next, validate the end-to-end
workflow with authorised reviewers using synthetic or otherwise
approved material. Then conduct a limited pilot with documented
oversight, without allowing scores alone to drive employment action.
Review performance, queue burden, false-positive burden under the
chosen label, data quality, and reviewer feedback. Only after those
stages should an organisation consider broader use, and it should
retain the ability to suspend the system if inputs, outputs, or
controls no longer meet expectations.

The prototype is therefore best understood as a research and
educational demonstration. It illustrates a multi-source feature
pipeline, group-aware model evaluation, a local service, and a review
interface. It does not establish that employee activity can be
reliably inferred from counts, that the dashboard rule is calibrated,
or that an organisation is prepared to act on its output. A responsible
next step is further evaluation under a defined governance framework,
not unreviewed deployment.

\section{Model and evaluation improvements}
Model work should first establish stronger evaluation discipline.
Experiments should compare the current random grouped split with
time-aware and release-aware designs, preserve an untouched test set,
and report uncertainty intervals where the sampling design supports
them. Synthetic augmentation should be compared with unaugmented
baselines through repeated validation, with data-generation
assumptions documented. Performance should be reported at several
review capacities and accompanied by calibration analysis only if a
probability interpretation is intended.

The dashboard's employee-level consensus rule requires its own
evaluation. Candidate thresholds should be chosen on validation data,
analyst capacity should be specified before test evaluation, and
employee-level labels should be aggregated under an explicit rule.
Evaluation should also examine how queue membership changes between
batches and whether changes are caused by data composition rather
than model behaviour.

\section{Software and governance improvements}
Engineering priorities include a locked environment, automated
regression tests, structured logs, configurable input limits, secure
temporary storage, consistent error handling, and a versioned
model-bundle manifest. The web service should receive a threat model
before any network exposure. A production workflow should add
authentication, role permissions, protected exports, durable audit
events, and a documented retention schedule.

Governance work should define the permitted purpose, data sources,
access rights, legal basis, review standards, appeal or correction
process, and criteria for ending a review. Any use involving real
people should undergo appropriate institutional and legal review,
impact assessment, and consultation with relevant stakeholders. A
human reviewer must be able to reject a model lead and record a
reasoned disposition.
